\subsection{R-34: progression-free Weyl pencils give a superlinear stable row range}

\paragraph{Statement R-34.}
Let $m\geq3$. Suppose $S\subset(\mathbb Z/m\mathbb Z)^2$ has
$r\geq3$ elements, contains $(0,0),(1,0),(0,1)$, and satisfies
\begin{equation}\label{R34:AP}
 2s=t+u,\qquad s,t,u\in S
                   \quad\Longrightarrow\quad s=t=u.
\end{equation}
The endpoints $t,u$ are allowed to be equal in this condition.
There are explicit cyclotomic size-$m$ pencils in $r$ variables
whose permanent and determinant forms are irreducible and have
finite complex projective stabilizer. For either form, positive
diagonal variable rescaling realizes every positive real full
moment vector with coordinate sum $m$.

Consequently the following conclusions hold, with separate
permanent and determinant witnesses.
\begin{enumerate}
\item For $n>m$, $a=n-m$, and every $\nu\vdash ns$ with
$r$ positive parts and $\nu_1>as$, the padded permanent image
and the degree-$n$ determinant image have stable enhanced
points with no nontrivial scalar multiple in their respective
$\mathrm{SL}_r$ orbits. Both multiplicities occur at every
integer stretch $j\geq\beta^2$ for a common ambient
invariant-generator bound $\beta$.
\item For every $\nu\vdash ms$ with $r$ positive parts,
both unpadded enhanced images have stable points and hence
the exact stable quotient dimensions. If $\nu$ is
nonrectangular, they also have no-scaling stable points and
the same all-integer occurrence conclusion. For rectangles,
no degree-period claim is made.
\end{enumerate}
Here the letter $s$ in $\nu\vdash ns$ or $ms$ denotes the
positive outer degree, independently of a label $s\in S$.

Such a set exists for every
\begin{equation}\label{R34:K}
 3\leq r\leq K_m:=4^{\lfloor\log_3m\rfloor}.
\end{equation}
Writing $\alpha=\log_3 4>1$, one has
$m^\alpha/4\leq K_m\leq m^\alpha<m^2$.
In particular $K_m>2m$ for every $m\geq2187$, and already
$K_{27}=64>54$. The theorem therefore extends the proved
strict-direction and unpadded nonrectangular witness ranges
beyond a linear number of rows. It does not assert a common
polynomial point for the two cones.

\begin{proof}
We first construct the pencils, then prove stabilizer and
moment assertions directly.

\paragraph{1. Clock-shift symmetry and the polynomial support.}
Let $\zeta$ be a primitive $m$th root of unity. On a basis
$e_i$, indexed modulo $m$, put
\[
 Xe_i=e_{i+1},\qquad Ze_i=\zeta^i e_i,
              \qquad U_{(a,b)}=X^aZ^b.
\]
These are monomial unitary matrices. Use variables $x_t$ for
$t\in S$, and set
\[
 A_S(x)=\sum_{t\in S}x_tU_t,\qquad
 q_{\det}(x)=\det A_S(x),\qquad
 q_{\mathrm{per}}(x)=\operatorname{per}A_S(x).
\]
For $h=(u,v)$, conjugation by $X^uZ^v$ multiplies $U_{(a,b)}$
by $\zeta^{va-ub}$. Determinant is invariant under this
conjugation. Permanent is invariant as well: a monomial
conjugation consists of inverse row-column permutations
and diagonal scalings whose total row-column product is one.
Thus both forms are fixed by the diagonal variable action
\[
              x_{(a,b)}\longmapsto\zeta^{va-ub}x_{(a,b)}.
\]
The characters on the distinct labels of $S$ are distinct,
since this pairing on $(\mathbb Z/m\mathbb Z)^2$ is
nondegenerate, also for composite $m$.

Let $q$ denote either form. Its coefficient $c_t$ of $x_t^m$
is the determinant or permanent of the monomial matrix $U_t$,
and is nonzero. Every monomial $x^\gamma$ in its support has
\begin{equation}\label{R34:weight_zero}
                    \sum_{t\in S}\gamma_t t=0
                 \quad\hbox{in }(\mathbb Z/m\mathbb Z)^2.
\end{equation}
In particular $x_s^{m-1}x_t$ is absent when $s\ne t$.

\paragraph{2. The progression-free condition forces finite projective stabilizer.}
Suppose $Bq=cq$ for a linear infinitesimal variable operator
\[
                  B=\sum_{t,u\in S}B_{tu}x_t\partial_{x_u}.
\]
Fix $s\ne t$ and compare the coefficient of
$x_s^{m-1}x_t$. The pure power contributes $m c_sB_{ts}$.
The other possible source monomials have the form
$x_s^{m-2}x_t x_u$, including $u=t$, or the already absent
$x_s^{m-1}x_u$ with $u\ne s$.
By \eqref{R34:weight_zero}, a monomial of the first type
would require
\[
                         2s=t+u.
\]
Condition \eqref{R34:AP} excludes it. The target monomial
itself has zero coefficient in $cq$, so
$m c_sB_{ts}=0$. Hence every off-diagonal entry of $B$ is zero.
Comparison of the nonzero pure-power coefficients then gives
$mB_{tt}=c$ for every $t$. Thus $B$ is scalar.

The complex projective algebraic stabilizer has zero Lie
algebra, so it is finite. This proof includes the possible
order-two difference when $m$ is even: the case $u=t$ was
explicitly included in \eqref{R34:AP}. It does not assume
that $m$ is prime.

\paragraph{3. The same pencils are irreducible.}
Keep only the variables with labels $(0,0),(1,0),(0,1)$,
calling them $x,y,z$. Their pencil is $xI+yX+zZ$.
Only the diagonal and the full cycle are perfect matchings,
so the two specialized forms are
\[
 \begin{aligned}
 q_{\mathrm{per}}(x,y,z)&=x^m-(-z)^m+y^m,\\
 q_{\det}(x,y,z)&=x^m-(-z)^m+(-1)^{m-1}y^m.
 \end{aligned}
\]
Both are irreducible. Indeed the constant term as a polynomial
in $x$ is a squarefree product of linear forms in $y,z$;
Eisenstein at any one of those linear primes applies.
If the full homogeneous $q$ had two positive-degree factors,
their restrictions to this three-variable subspace would be
nonzero homogeneous factors of the same positive degrees:
neither can vanish because their product is nonzero.
This contradicts the displayed irreducibility. Hence the
full forms are irreducible. Invertible diagonal variable
rescaling preserves this property and the finite stabilizer.

\paragraph{4. A complete convex proof of prescribed moments.}
Use the tensor-induced Hermitian metric on
$\operatorname{Sym}^m\mathbb C^r$. The finite diagonal
symmetry of paragraph~1 is unitary and has distinct characters.
Moment equivariance therefore makes $M(q)$ diagonal. The same
is true after every positive diagonal variable rescaling,
since that rescaling commutes with the symmetry.

Write $q=\sum_\gamma c_\gamma x^\gamma$ and put
$w_\gamma=|c_\gamma|^2\|x^\gamma\|^2>0$ on its support.
For $t\in\mathbb R^r$ with $\sum_i t_i=0$, let
$q_t(x)=q(e^{t_1}x_1,\ldots,e^{t_r}x_r)$. Then
\[
 \|q_t\|^2=\sum_\gamma w_\gamma e^{2\langle\gamma,t\rangle},
 \qquad
 M(q_t)_{ii}
   =\frac{\sum_\gamma\gamma_iw_\gamma
                    e^{2\langle\gamma,t\rangle}}
          {\sum_\gamma w_\gamma e^{2\langle\gamma,t\rangle}}.
\]
Let $\mu_i>0$ have sum $m$, and minimize on that hyperplane
\begin{equation}\label{R34:convex}
 \Phi(t)=\log\!\left(\sum_\gamma w_\gamma
                         e^{2\langle\gamma,t\rangle}\right)
                                      -2\langle\mu,t\rangle.
\end{equation}
Its Hessian is four times the covariance matrix of the
exponent distribution. All pure exponents $m e_i$ have
positive weight. Therefore its quadratic form is strictly
positive on every nonzero vector in $\sum_i t_i=0$:
a zero variance would make all $m t_i$ equal, hence $t=0$.
Thus $\Phi$ is strictly convex there.

It is also coercive. Let $M_t=\max_i t_i$ and
$w_* =\min_i w_{m e_i}>0$. Then
\[
 \Phi(t)\geq\log w_*+2\bigl(mM_t-\langle\mu,t\rangle\bigr)
       =\log w_*+2\sum_i\mu_i(M_t-t_i).
\]
For nonzero $t$ in the zero-sum hyperplane the last sum is
strictly positive. As a positive continuous homogeneous
function on that hyperplane, it has a positive minimum on
the unit sphere. The displayed bound therefore tends to
infinity with $\|t\|$. A unique minimizer exists.
At it the projected gradient is zero. The coordinate sum
of $M(q_t)-\mu$ is zero, so that projected equation is the
full equality
\begin{equation}\label{R34:prescribed}
                           M(q_t)=\operatorname{diag}(\mu).
\end{equation}
This proves the moment assertion for every positive real
$\mu$, without an unstated moment-polytope surjectivity theorem.
Normalize $q_t$ to norm one if desired; its moment is unchanged.

\paragraph{5. Separate stable witnesses and elimination of their periods.}
For the padded assertion, index variables so the first
coordinate corresponds to the first part of $\nu$, and take
\[
 \mu_1=\nu_1/s-a,\qquad \mu_i=\nu_i/s\quad(i>1).
\]
These entries are positive and sum to $m$. Apply paragraph~4
separately to $q_{\mathrm{per}}$ and $q_{\det}$, obtaining
irreducible forms of finite complex projective stabilizer
with that same full moment matrix.

For either resulting $q$, put $\ell=x_1$ and take the standard
unit auxiliary extremal vector $v_\eta$, where
$\eta=\nu-\nu_r\mathbf1_r$. The ordered tensor
\[
          (\ell^{\otimes a}\otimes q)^{\otimes s}\otimes v_\eta
\]
has full moment $\nu_r I_r$. Its projective stabilizer is
finite because it preserves $[q]$. Kempf--Ness gives stability.
The initial term of a polynomial product is the nonzero
product of the initial terms for every one-parameter subgroup.
Thus the same minimum-weight argument as G-22 transfers
stability to
\[
                      (\ell^a q)^{\otimes s}\otimes v_\eta.
\]
For the permanent this is in the padded permanent image.
For the determinant it is in the degree-$n$ determinant
image by block-diagonal padding with $a$ linear entries.

Fix $q$ in each case. Its finite projective automorphism
group acts on the padding lines. Stability is open in this
line parameter and holds at the line just constructed.
Avoid the finitely many proper fixed loci of nonidentity
automorphisms. Since $q$ is nonlinear and irreducible,
any projective symmetry of $\ell^a q$ preserves its unique
nonlinear factor and its padding line, even when $a=1$.
The resulting padded form is projectively free. An enhanced
tensor symmetry in $\mathrm{SL}_r$ is therefore central,
and the central exponent $ns-|\eta|=r\nu_r$ is divisible
by $r$. It acts trivially. This gives separate stable
no-scaling points, including when $\eta=0$.

For the unpadded assertion use instead $\mu_i=\nu_i/s$.
Then $q^{\otimes s}\otimes v_\eta$ has moment zero and
finite projective stabilizer directly. If $\eta\ne0$, fix
this $q$ and vary the auxiliary flag in its nonempty stable
open, avoiding the finite automorphism fixed loci as in
R-32. The center exponent is $ms-|\eta|=r\nu_r$, so the
same no-scaling conclusion follows. If $\eta=0$, retain
stability without making a period-one assertion.

In each of the no-scaling cases, separately in each ring,
the closed-orbit separation and numerical-semigroup argument
of R-28 gives occurrence at every $j\geq\beta^2$. No
invariant is required to be simultaneously nonzero on the
two different witnesses.

\paragraph{6. An explicit progression-free set for every modulus.}
Let $k=\lfloor\log_3 m\rfloor$, and let $A$ be the set of
nonnegative integers whose $k$ base-three digits lie in
$\{0,1\}$. Then
\[
 |A|=2^k,\qquad 0\leq A\leq(3^k-1)/2\leq(m-1)/2.
\]
If $x+z\equiv2y\pmod m$ for $x,y,z\in A$, the two sides
differ by an integer of absolute value less than $m$, so
they are equal as integers. There are no carries in either
$x+z$ or $2y$ in base three. Equality of every digit forces
$x=z=y$. Hence $A$ is progression-free modulo $m$ in the
strong sense including equal endpoints.
The product $A\times A$ has the same property, size $4^k$,
and contains $(0,0),(1,0),(0,1)$.
Any subset retaining those three labels gives every size
in \eqref{R34:K}.

Finally $3^k\leq m<3^{k+1}$ gives the displayed power bounds.
For $k\geq7$, $(4/3)^k>6$, so
$4^k>2\cdot3^{k+1}>2m$. This proves the uniform threshold
$m\geq3^7=2187$ and the explicit superlinear row count.
\end{proof}

\paragraph{Scope and source use.}
Clock-shift conjugation, the coefficient comparison, the
convex normalization, and the digit construction were proved
explicitly above. The only cited invariant-theory inputs are
the already registered Kempf--Ness, Hilbert--Mumford,
stability-openness, and closed-orbit separation results
(G-B14/G-B15), together with the already proved semigroup
argument. No classical additive-combinatorial density bound
is imported. The pencils before normalization have explicit
cyclotomic coefficients; the minimizer is an existence
construction, not a polynomial bit-complexity claim.
The result gives stable occurrence witnesses and makes no
claim of a cancellation-free coefficient formula or a
complexity-class separation.
